\chapter{Run 3 calibration of the strip \texorpdfstring{d$E$/d$x$}{dE/dx}}
\label{app:run3_calib}

The harmonic estimator \ih\ relies on a calibrated charge per unit path length in each
silicon module. In Run 2 this charge was stable over the whole data-taking period, so no
time-dependent correction was needed. This is no longer the case in Run 3, and a dedicated
calibration was derived for the 2022--2024 data. The estimator is defined as
\begin{equation}
    I_{h}= \biggl( \cfrac{1}{N} \sum_j^{N} \Bigl(\cfrac{\mathrm{d}E}{\mathrm{d}x}(j)\Bigr)^{-2} \biggr)^{-1/2},
    \label{eq:app_harmonic}
\end{equation}
where $\mathrm{d}E/\mathrm{d}x(j)$ is the calibrated charge per unit path length of the
$j^{\text{th}}$ measurement along a track.

Figure~\ref{fig:app_stab_nocalib} compares the stability of \ih\ between the two runs. The
\ih\ mean is extracted from a Gaussian fit, plotted against the run number for Run 2
(run number $<310$k for 2017, $>310$k for 2018) and against the integrated luminosity for
Run 3. In Run 2 the variation stays around 0.3\%, while Run 3 drifts downwards by several
percent. The difference comes from the reconstruction: the Run 2 samples use the
Ultra-Legacy reprocessing, whereas the Run 3 samples are produced with prompt reconstruction.

\begin{figure}[h!]
    \centering
     \includegraphics[width=0.47\textwidth]{figures/Run3_StripCalibratrion/Stab_Ih_Gauss_Run2.pdf}
     \includegraphics[width=0.47\textwidth]{figures/Run3_StripCalibratrion/Stab_IhNoCalib_Run3_Lumi.pdf}
    \caption{Stability of \ih\ before calibration: (left) Run 2~\cite{CMS-AN-22-129}, from
    2017--2018 Single Muon UL AOD samples with $\pt>55\GeV$ and $|\eta|<1$; (right) Run 3,
    from 2022--2024 ZMu skim Prompt Reco AOD samples with $25<\pt<55\GeV$ and $|\eta|<2.4$.}
    \label{fig:app_stab_nocalib}
\end{figure}

The d$E$/d$x$ information is only available in the AOD format. Following the current rules on
data-format availability, the calibration is derived from the ZMu skim, which gives access to
muon tracks with lower \pt than the analysis selection, at the price of a smaller sample than
the Run 2 Single Muon datasets. The \ih\ stability uses tracks with $25<\pt<55\GeV$,
$|\eta|<2.4$, High Purity, $N_{\text{Pix}}\ge 2$, $N_{\text{validhits}}\ge 8$,
$\chi^2/\text{ndf}\le 5$, $|d_{xy}|\le 0.02$, $|d_z|\le 0.1$, $E/P\le 0.3$,
MiniRel Iso $\le 0.1$, $\sum_{\Delta R<0.3}\pt\le 15\GeV$ and $\sigma(\pt)/\pt\le 0.25$.
For the per-layer calibration the selection is loosened to gain statistics:
$15<\pt<55\GeV$, $|d_{xy}|\le 0.5$, $|d_z|\le 0.5$, and the $E/P$ and isolation-sum
requirements are dropped.

The luminosity binning is chosen to contain enough events per bin for the fit, while staying
narrow enough to follow changes such as new gain settings or HV adjustments. It is shown in
Fig.~\ref{fig:app_lumibin} and corresponds to bins of 3--5\fbinv.

\begin{figure}[h!]
    \centering
     \includegraphics[width=0.55\textwidth]{figures/Run3_StripCalibratrion/lumi_ranges.pdf}
    \caption{Luminosity binning (vertical lines) used to calibrate the d$E$/d$x$ per strip
    layer over the 2022--2024 period.}
    \label{fig:app_lumibin}
\end{figure}

In each luminosity bin, the peak of the d$E$/d$x$ distribution is fitted with a Gaussian, and
its mean is taken as the reference for that layer. The corresponding MPV values are shown in
Fig.~\ref{fig:app_mpv} for the TIB and TOB layers and for the TID and TEC rings. The pixel
detector is not calibrated and is not shown here.

\begin{figure}[h!]
    \centering
     \includegraphics[width=0.47\textwidth]{figures/Run3_StripCalibratrion/ChargTib_zoom_Lumi.pdf}
     \includegraphics[width=0.47\textwidth]{figures/Run3_StripCalibratrion/ChargTob_zoom_Lumi.pdf}\\
     \includegraphics[width=0.47\textwidth]{figures/Run3_StripCalibratrion/ChargRTid_zoom_Lumi.pdf}
     \includegraphics[width=0.47\textwidth]{figures/Run3_StripCalibratrion/ChargRTec_zoom_Lumi.pdf}
    \caption{Gaussian peak of the d$E$/d$x$ distribution as a function of the integrated
    luminosity: TIB layers (top left), TOB layers (top right), TID rings (bottom left) and
    TEC rings (bottom right). Each layer or ring is drawn in a separate colour.}
    \label{fig:app_mpv}
\end{figure}

A correction factor is defined in each bin as $\text{CF}=3.10/\mu_{\text{Gauss}}$, where the
value 3.10 sets the normalisation. Once these factors are applied, the Run 3 \ih\ stability
improves markedly and becomes comparable to Run 2, as shown in
Fig.~\ref{fig:app_stab_calib}. The residual structure in the corrected Run 3 plot follows
from the different binnings: the correction factors are derived in bins of 3--5\fbinv, while
the stability is displayed with a 1\fbinv\ binning.

\begin{figure}[h!]
    \centering
     \includegraphics[width=0.47\textwidth]{figures/Run3_StripCalibratrion/Stab_Ih_Gauss_Run2.pdf}
     \includegraphics[width=0.47\textwidth]{figures/Run3_StripCalibratrion/Stab_Ih_Run3_Lumi.pdf}
    \caption{Stability of \ih\ after calibration: (left) Run 2, repeated from
    Fig.~\ref{fig:app_stab_nocalib} for convenience~\cite{CMS-AN-22-129}; (right) Run 3 with
    the correction factors applied.}
    \label{fig:app_stab_calib}
\end{figure}

This calibration was presented in the ``Strip calibration and local reconstruction''
meetings of the Tracker DPG, the last time on March 25, 2025
(\href{https://indico.cern.ch/event/1505992/#b-603415-strip-calibration-loc}{indico agenda}).
The Tracker DPG convener, Martin Delcourt, gave his green light on the proposed correction at
the end of that meeting.

